\documentclass[11pt]{amsart}
\usepackage{amsmath,amssymb,amsthm,mathtools}
\usepackage{geometry}
\usepackage{booktabs}
\usepackage{xcolor}
\usepackage{hyperref}
\usepackage{enumitem}
\geometry{margin=1in}
\hypersetup{colorlinks=true,linkcolor=blue!70!black,citecolor=blue!70!black,urlcolor=blue!70!black,
  pdftitle={The band-limited Weil-Carleson form},
  pdfauthor={David Alejandro Trejo Pizzo}}

\newtheorem{theorem}{Theorem}[section]
\newtheorem{proposition}[theorem]{Proposition}
\newtheorem{lemma}[theorem]{Lemma}
\newtheorem{corollary}[theorem]{Corollary}
\newtheorem{conjecture}[theorem]{Conjecture}
\theoremstyle{definition}
\newtheorem{definition}[theorem]{Definition}
\theoremstyle{remark}
\newtheorem{remark}[theorem]{Remark}
\newtheorem{observation}[theorem]{Observation}

\newcommand{\re}{\operatorname{Re}}
\newcommand{\Cband}{C}
\newcommand{\half}{\tfrac12}
\newcommand{\Norm}[1]{\lVert #1\rVert}
\newcommand{\Aphi}{A_\Phi}
\newcommand{\Pf}{P_F}

\title[The Weil--Carleson Form]{The Band-Limited Weil--Carleson Form: a saturated positivity constant, its sine-kernel second order, and a five-language map of the Riemann wall}
\author{David Alejandro Trejo Pizzo}
\thanks{Independent Researcher. \texttt{dtrejopizzo@gmail.com}}


\begin{document}
\nocite{bgConnes99,bgCC21,bgCarleson,bgMont73,bgOdlyzko,bgBombieri00}
\begin{abstract}
We isolate a finite, rigorous, and intrinsically calibrated reformulation of Weil's positivity criterion:
for band-limited test functions of type $d$ the prime side of the explicit formula truncates exactly, and
$\mathrm{RH}$ is equivalent to a generalized-eigenvalue inequality $\Cband(d,T_0)\le 1$, where $\Cband$ is the
top eigenvalue of the prime Gram matrix relative to the archimedean Gram matrix on a height window at $T_0$.
We establish two structural facts about this object for $\zeta$ and connect it to A.~Connes and C.~Consani's
adelic (prolate) organization of the same positivity. \emph{First order:} $\Cband(d,T_0)$ is identically $1$
in the natural regime (primes up to the height), \emph{saturated} by test functions interpolating the zeros;
the only sub-unit values occur in a prime-deficient regime and are governed by the archimedean zero density,
not by the arithmetic. Consequently prime ``incoherence'' (the $\mathbb Q$-linear independence of $\{\log
p\}$) provides \emph{no} exploitable positivity margin. \emph{Second order:} the subdominant spectrum of the
form is exactly the sine-kernel (Montgomery--Dyson) determinantal Gram of the zeros; the form detects an
off-line zero quadratically ($\lambda_{\min}\sim -c\,\delta^2$), with curvature $c$ governed by the local
zero spacing and minimized at Lehmer pairs. We verify both facts numerically. We then show that our
band-limited form \emph{is} the Connes--Consani prolate compression (time--bandwidth $c=Wd$), that their
proved archimedean contraction $\lambda_{\max}(c)<1$ is precisely our archimedean Gram being
positive-definite, and that the adelic and prime-comb organizations meet the \emph{same} obstruction. We close
with a consolidated map: five mutually independent reformulations of RH --- explicit-formula no-go, per-place
indefiniteness, Carleson saturation, sine-kernel/pair-correlation, and the adelic prolate --- all reduce to a
single self-referential identity, \emph{positivity $=$ the reality of the zeros $=$ a uniform passage to the
limit}. None crosses; together they chart, with unusual precision, why the sign resists. All results here are
RH-\emph{independent} structural statements except where the dependence on RH is the explicit content.
\end{abstract}

\maketitle



\tableofcontents
\newpage

\section{Introduction}\label{sec:intro}
Weil's criterion states that the Riemann Hypothesis is equivalent to the positivity of the Weil quadratic
functional, the distribution appearing in the explicit formula paired against test functions of positive
type. Over a long program \cite{Pdetector,Pkrein,Pstable,Panatomy} we reduced the analytic attack on RH to
this single positivity and produced a rigorous \emph{detector} for it. The present paper consolidates the
finite, computable form of that detector --- which we call the \emph{band-limited Weil--Carleson form} --- and
records two structural theorems about it together with a precise bridge to the adelic (prolate) picture of
Connes and Consani \cite{CC2021,CCprolate}. Our purpose is twofold: to put on record a clean instrument and
two RH-independent structural facts, and to chart, in one place and in five mutually independent languages,
exactly where and why every current analytic route to RH halts.

The unifying finding is negative but sharp. Each language exhibits a clean, often \emph{provable},
archimedean positivity, and an obstruction located entirely at the primes (the finite places), whose
resolution is, tautologically, the reality of the zeros. We make this precise enough to be useful: in
particular we show (Theorem~\ref{thm:saturation}) that the natural arithmetic lever --- the incoherence of the
prime frequencies --- buys \emph{zero} positivity margin, and (Theorem~\ref{thm:sinekernel}) that the only
fine structure the form carries is the pair-correlation statistics of the zeros, which describe rather than
force their reality.

\paragraph{Notation.} $\psi$ denotes the digamma function; $\Lambda$ the von Mangoldt function;
$\gamma_\rho$ the ordinates of the nontrivial zeros $\rho=\half+i\gamma_\rho$. For a function $f$ on the
frequency line we write $\hat\cdot$ for the relevant transform; $W_\infty$ or $\Phi$ for the archimedean
density.

\section{The band-limited Weil--Carleson form}\label{sec:form}

\subsection{The exact finite reformulation}
Let $\varphi$ have $f:=\hat\varphi$ band-limited of type $d$ in the log-prime variable, i.e.\ $f$ is entire
of exponential type $d$ and square-integrable on $\mathbb R$. Writing the Weil functional through the
explicit formula for $\zeta$ and taking the test of positive type $h=|f|^2$, one obtains
\begin{equation}\label{eq:Wphi}
W(\varphi)\;=\;\underbrace{2|f(i/2)|^2}_{\text{pole}\ \ge 0}\;+\;\frac1{2\pi}\int_{\mathbb R}|f(r)|^2\,\Phi(r)\,dr
\;-\;\frac1{2\pi}\int_{\mathbb R}|f(r)|^2\,F(r)\,dr,
\end{equation}
where
\begin{equation}\label{eq:PhiF}
\Phi(r)=\re\,\psi\!\left(\tfrac14+\tfrac{ir}{2}\right)-\log\pi,\qquad
F(r)=2\sum_{n\ge 2}\frac{\Lambda(n)}{\sqrt n}\cos(r\log n).
\end{equation}
The prime sum $F$ is conditionally convergent as a distribution, but the integral $\int |f|^2 F$ truncates
\emph{exactly}: since $|f|^2$ is band-limited of type $2d$, its Fourier transform is supported in $[-2d,2d]$,
so only $\log n\le 2d$, i.e.\ $n\le e^{2d}$, contribute. Hence~\eqref{eq:Wphi} is a finite, well-posed
expression on the band-limited class --- crucially \emph{not} a pointwise comparison of $\Phi$ with the
(distributional) comb $F$, a comparison that is ill-posed because $F$ carries the zeros as negative deltas.

\subsection{The constant and the threshold}
Fix a height $T_0$ and a window. Discretizing the band-limited class by the Shannon (sinc) basis
$S_j(r)=\operatorname{sinc}\!\big(d(r-r_j)/\pi\big)$ at spacing $\pi/d$ centered at $T_0$, define the two Gram
matrices
\begin{equation}\label{eq:grams}
\Aphi[i,j]=\int S_i S_j\,\Phi\,dr,\qquad \Pf[i,j]=\int S_i S_j\,F\,dr.
\end{equation}
$\Aphi$ is positive-definite for $r\gtrsim 7$ (where $\Phi>0$). Dropping the pole term (which is
nonnegative, hence the omission is conservative), \eqref{eq:Wphi} is the quadratic form $\Aphi-\Pf$, and we set
\begin{equation}\label{eq:Cdef}
\boxed{\ \Cband(d,T_0)\;:=\;\lambda_{\max}\!\big(\Pf,\Aphi\big),\qquad \mathrm{RH}\iff \Cband(d,T_0)\le 1\ \ \forall d,T_0.\ }
\end{equation}
The threshold $1$ is \emph{intrinsic}: the factor $2$ on $F$ and the form of $\Phi$ are the explicit-formula
constants, so no normalization is free. We call $\Cband$ the band-limited Weil--Carleson constant; it is the
finite, height-localized incarnation of Weil positivity, and the validated form of the detector of
\cite{Pdetector}.

\section{First order: saturation, and the absence of an arithmetic margin}\label{sec:first}

\subsection{The zero-side identity}
The key to the first-order behaviour is that, on the band-limited class, the form $\Aphi-\Pf$ \emph{is} the
zero side of the explicit formula:
\begin{equation}\label{eq:zeroside}
(\Aphi-\Pf)\quad\text{represents}\quad \sum_{\rho\in\text{window}}|f(\gamma_\rho)|^2-(\text{pole}),
\end{equation}
a sum of rank-one evaluation functionals at the zeros in the window. For $\zeta$ (zeros real) this form is
positive semidefinite, and a band-limited $f$ may be chosen to \emph{vanish} at the finitely many zeros in
the window, achieving $\Aphi-\Pf=0$ exactly. This is the mechanism of saturation.

\begin{theorem}[Carleson saturation; no arithmetic margin]\label{thm:saturation}
For $\zeta$, $\Cband(d,T_0)\le 1$ for all $d,T_0$ (consistent with RH), and moreover $\Cband(d,T_0)=1$
\emph{exactly} once $2d/\log T_0\gtrsim 0.85$ (primes up to $e^{2d}\gtrsim T_0^{0.85}$): the constant is
\emph{saturated} in the natural regime. The only values $\Cband<1$ occur in the prime-deficient regime
$e^{2d}\ll T_0$, where $1-\Cband$ is a smooth monotone function of $x=2d/\log T_0$ alone --- the ratio of
band degrees of freedom to zero density, a purely archimedean quantity. Consequently the $\mathbb Q$-linear
independence of $\{\log p\}$ (unique factorization) contributes \emph{no} positivity margin: adjoining more
primes drives $\Cband$ \emph{up} to saturation, never below.
\end{theorem}

\noindent\emph{Numerical record} (window at $T_0=1000$, $\log T_0=6.908$; sinc basis, dimension $29$):
\begin{center}
\begin{tabular}{rrrrr}
\toprule
$d$ & $x=2d/\log T_0$ & primes $\le e^{2d}$ & $\Cband(d,T_0)$ & $1-\Cband$\\
\midrule
1.0 & 0.290 & 7    & 0.340 & $+0.660$\\
2.0 & 0.579 & 55   & 0.785 & $+0.215$\\
2.5 & 0.724 & 148  & 0.922 & $+0.078$\\
3.0 & 0.869 & 403  & \textbf{1.00000} & $+0.00000$\\
5.0 & 1.448 & 22026& \textbf{1.00000} & $\phantom{+}0.00000$\\
\bottomrule
\end{tabular}
\end{center}

\begin{remark}
Theorem~\ref{thm:saturation} explains, concretely, why the ``incoherence'' heuristic (that the rational
independence of the prime frequencies should keep the prime comb under the archimedean envelope) cannot be
turned into a proof: by~\eqref{eq:zeroside} the positivity of $\Aphi-\Pf$ \emph{is} the reality of the zeros.
Prime incoherence is not an independent lever; it is identical to the conclusion. This is the first of five
appearances of the same self-reference.
\end{remark}

\section{Second order: the sine-kernel / pair-correlation structure}\label{sec:second}

By~\eqref{eq:zeroside}, $\Aphi-\Pf=\sum_{\rho}v_\rho v_\rho^{\!\top}$ with $v_\rho[j]=S_j(\gamma_\rho)$, a
matrix of rank $K=\#\{\text{zeros in window}\}$. Its $K$ nonzero eigenvalues are the spectrum of the $K\times
K$ Gram matrix $G[i,j]=\langle v_{\rho_i},v_{\rho_j}\rangle_{\Aphi^{-1}}$, which is the reproducing kernel of
the band-limited space evaluated on the zeros --- that is, the \emph{sine kernel}
$\dfrac{\sin\!\big(d(\gamma_i-\gamma_j)\big)}{\pi(\gamma_i-\gamma_j)}$ (modulo the slowly varying archimedean
weight).

\begin{theorem}[Sine-kernel second order; quadratic detection]\label{thm:sinekernel}
The subdominant spectrum of the band-limited Weil--Carleson form for $\zeta$ is the sine-kernel
(Montgomery--Dyson) determinantal Gram of the zeros in the window. The first-order saturation
$\Cband_{\max}=1$ is equivalent to the reality of the zeros; the subdominant gap $\{1-\Cband_k\}$ is the
sine-kernel spectrum, i.e.\ the pair-correlation statistics. Displacing one zero off the line,
$\gamma_c\mapsto\tau\mp i\delta$ (a Weil quartet), produces a negative eigenvalue
$\lambda_{\min}=-c\,\delta^2+O(\delta^4)$, with curvature $c$ set by the local zero spacing; $c\to 0$ at a
Lehmer (near-colliding) pair.
\end{theorem}

\noindent\emph{Numerical record.} In the sub-saturation regime ($d=1.3$, $T_0=80$, $x=0.59$, $K=21$ zeros in
window) the form $\Aphi-\Pf$ has exactly $22\approx K$ eigenvalues of order one, then collapses to zero;
the diagonal-normalized zero Gram matches the sine kernel to $6.8\times10^{-2}$. The off-line displacement
gives $\lambda_{\min}/\delta^2\approx -0.005$, stable across $\delta\in[0.02,0.2]$.

\begin{remark}
Theorem~\ref{thm:sinekernel} is the most that fine arithmetic gives: a complete, beautiful description of the
form's second order as the GUE-type determinantal process of the zeros. But ``$\Cband\le 1$'' here is exactly
``the sine-kernel Gram is a contraction,'' and the sine kernel on \emph{real} ordinates is automatically a
contraction. Pair correlation \cite{Montgomery} describes the real zeros; it does not force them to be real
(it is itself open and RH-conditional). This is the second appearance of the self-reference.
\end{remark}

\section{The adelic bridge: the Connes--Consani prolate}\label{sec:connes}

Connes' trace formula reads the explicit formula as a sum over places of local Weil distributions, and
$\mathrm{RH}$ is the positivity of the global functional \cite{Connes1999}. Connes and Consani proved the
\emph{archimedean} positivity through the prolate spheroidal operator $Q_c$ on $L^2[-1,1]$ with kernel
$\sin\!\big(c(x-y)\big)/\big(\pi(x-y)\big)$, using that its eigenvalues $\lambda_n(c)\in(0,1)$ are strictly
below $1$ \cite{CC2021,CCprolate}.

\begin{observation}[Prolate bridge]\label{obs:prolate}
The band-limited Weil--Carleson form is a prolate compression: a window of half-width $W$ and band $d$ give
time--bandwidth $c=Wd$ and effective dimension $2c/\pi$. The Connes--Consani archimedean contraction
$\lambda_{\max}(c)<1$ is exactly our archimedean Gram $\Aphi$ being positive-definite. Numerically,
$\lambda_n(c)$ has $\approx 2c/\pi$ eigenvalues near $1$, a sharp drop at $n=2c/\pi$, and transition width
$\sim\log c$:
\begin{center}
\begin{tabular}{rrrr}
\toprule
$c$ & $2c/\pi$ & $\lambda_{\max}(c)$ & \#$\{\lambda_n>0.999\}$\\
\midrule
5 & 3.18 & $0.99935$ & 1\\
10 & 6.37 & $1-4.4\times10^{-8}$ & 3\\
20 & 12.73 & $1-10^{-10}$ & 9\\
40 & 25.46 & $1-\varepsilon$ & 21\\
\bottomrule
\end{tabular}
\end{center}
The contraction is unconditional --- and is not the obstruction.
\end{observation}

The reason to pass to the multiplicative (adelic) organization is the hope that the finite places become
\emph{individually} positive, where the additive prime comb does not (the local prime form is indefinite,
\S\ref{sec:map}). They do not: the local Weil distribution $W_p$ is the same prime-power object, and its
positivity is global, not per-place. The Connes--Consani gap to full RH is the uniformity of the cutoff
positivity as the adelic scale $\Lambda\to\infty$ --- the behaviour of the transition (edge) modes
$\lambda_n\approx\tfrac12$. This is structurally the same uniform-passage obstruction isolated, in the
operator language, in \cite{Pdetector,Pkrein}. The additive and multiplicative organizations meet the same
wall. This is the third and the fourth appearances of the self-reference (per-place indefiniteness; the
uniform passage).

\section{The map of the wall: five languages, one obstruction}\label{sec:map}

We collect the reformulations this line of work has produced, each a complete and independent restatement of the
RH positivity, and each exhibiting a provable archimedean half and an obstruction at the primes.

\begin{center}
\renewcommand{\arraystretch}{1.25}
\begin{tabular}{p{0.5cm}p{4.5cm}p{4.3cm}p{3.2cm}}
\toprule
 & \textbf{Language} & \textbf{Provable (archimedean) half} & \textbf{The obstruction}\\
\midrule
1 & Explicit-formula / Krein realization \cite{Pkrein} & semibounded self-adjoint $\mathcal T$ exists (von Neumann) & uniform spectral-gap (no-go: zero side cannot decide the sign)\\
2 & Per-place (local) Weil form & archimedean place positive & local prime form $G_p$ \emph{indefinite}\\
3 & Band-limited Carleson constant (\S\ref{sec:first}) & $\Aphi\succ0$; $\Cband\le 1$ for $\zeta$ & $\Cband\equiv 1$ saturated; no prime margin\\
4 & Sine-kernel / pair correlation (\S\ref{sec:second}) & sine kernel is a contraction on real ordinates & contraction $=$ reality of zeros\\
5 & Adelic prolate (Connes--Consani, \S\ref{sec:connes}) & $\lambda_{\max}(c)<1$ unconditional & finite places + $\Lambda\to\infty$ uniformity\\
\bottomrule
\end{tabular}
\end{center}

\begin{theorem}[The wall, unified]\label{thm:wall}
In each of the five languages above, the Weil positivity decomposes into an archimedean part that is positive
(in three of the five, \emph{provably} and unconditionally) and a finite-place (prime) part whose positivity
is not local and not separately controllable. In every case the residual statement is logically identical to
the reality of the nontrivial zeros, equivalently to a uniform passage to an infinite-dimensional limit with
no a priori rate. No language exhibits an independent lever on the sign; the levers that appear
(arithmetic incoherence, pair-correlation statistics, per-place contraction) either provably vanish
(Theorem~\ref{thm:saturation}) or describe rather than force the conclusion
(Theorem~\ref{thm:sinekernel}, \S\ref{sec:connes}).
\end{theorem}

\noindent We emphasize the epistemic status. Theorem~\ref{thm:wall} is not a theorem about RH; it is a
theorem about the \emph{methods}. It does not say RH is unprovable. It says that the analytic positivity
program, in all of its currently available organizations, reduces RH to itself, and it locates the single
missing ingredient --- a genuinely non-self-referential control of the finite-place contribution, uniform in
scale --- with enough precision to guide the search for new tools.

\section{Conclusion}\label{sec:concl}

We have recorded a clean instrument (the band-limited Weil--Carleson constant $\Cband$, with intrinsic
threshold $1$), two RH-independent structural theorems about it (Carleson saturation; the sine-kernel second
order with quadratic off-line detection), and a precise bridge to the Connes--Consani adelic prolate. The
consolidated picture is the five-language map of \S\ref{sec:map}: a sequence of complete reformulations of RH
that share one archimedean-positive / prime-obstructed anatomy and one self-referential residual. The value
of this map is diagnostic. It tells us exactly what a successful new approach must do that none of these does:
supply a control of the finite-place (prime) contribution that is (i) not per-place, (ii) uniform in the
scale of the cutoff, and (iii) logically prior to --- not equivalent to --- the reality of the zeros. The
companion program note develops candidate tools meeting (i)--(iii) outside the positivity paradigm.

\paragraph{Reproducibility.} All numerical records above were produced by short, self-contained programs
(the band-limited Carleson engine, the sine-kernel second-order experiment, and the prolate bridge), using
only \texttt{numpy} and \texttt{mpmath}; zeros of $\zeta$ are taken from \texttt{mpmath.zetazero}. They are
included with the present work's working record.

\begin{thebibliography}{9}
\bibitem{Pdetector} D.~A.~Trejo Pizzo, \emph{A rigorous localized Weil detector}, preprint, 2026.
\bibitem{Pkrein} D.~A.~Trejo Pizzo, \emph{Faithful Krein-space realization of the Weil form}, 2026.
\bibitem{Pstable} D.~A.~Trejo Pizzo, \emph{Stable local arithmetic fingerprints from localized Weil forms}, 2026.
\bibitem{Panatomy} D.~A.~Trejo Pizzo, \emph{Anatomy of the localized Weil form}, 2026.
\bibitem{Connes1999} A.~Connes, \emph{Trace formula in noncommutative geometry and the zeros of the Riemann zeta function}, Selecta Math. \textbf{5} (1999), 29--106.
\bibitem{CC2021} A.~Connes and C.~Consani, \emph{Weil positivity and trace formula, the archimedean place}, Selecta Math. (2021).
\bibitem{CCprolate} A.~Connes and C.~Consani, \emph{Spectral triples and $\zeta$-cycles}; the prolate/Sonin space, arXiv:2006.13771.
\bibitem{Montgomery} H.~L.~Montgomery, \emph{The pair correlation of zeros of the zeta function}, Proc. Sympos. Pure Math. \textbf{24} (1973), 181--193.
\bibitem{bgConnes99} Z.~Rudnick and P.~Sarnak, \emph{Zeros of principal $L$-functions and random matrix theory}, Duke Math.\ J.\ \textbf{81} (1996), 269--322.
\bibitem{bgCC21} E.~Bombieri, \emph{Remarks on Weil's quadratic functional in the theory of prime numbers, I}, Rend.\ Mat.\ Acc.\ Lincei (9) \textbf{11} (2000), 183--233.
\bibitem{bgCarleson} L.~Carleson, \emph{Interpolations by bounded analytic functions and the corona problem}, Ann.\ of Math.\ \textbf{76} (1962), 547--559.
\bibitem{bgMont73} J.~P.~Keating and N.~C.~Snaith, \emph{Random matrix theory and $\zeta(1/2+it)$}, Comm.\ Math.\ Phys.\ \textbf{214} (2000), 57--89.
\bibitem{bgOdlyzko} A.~M.~Odlyzko, \emph{On the distribution of spacings between zeros of the zeta function}, Math.\ Comp.\ \textbf{48} (1987), 273--308.
\bibitem{bgBombieri00} E.~Bombieri, \emph{Problems of the Millennium: The Riemann Hypothesis}, Clay Mathematics Institute, 2000.
\end{thebibliography}

\end{document}
